% Pilot on the Windows tasks. Included in the results (journal) or the appendix (conference).
Table~\ref{tab:pilot} reports the pilot on the \ResPilotTasks{} Windows tasks. All tasks completed and every
verdict was correct (accuracy \ResPilotAccuracy{}): the eight attacks were judged malicious and the three
benign look-alikes benign, so the false-positive and false-negative rates were both zero. Mean quality was
\ResPilotQuality{}, mean efficiency \ResPilotEfficiency{} and mean reward \ResPilotReward{}, with
\ResPilotCalls{} tool calls per task (range \ResPilotMinCalls{}--\ResPilotMaxCalls{}). The whole run took
\ResPilotWall\,s of wall-clock time with four tasks in parallel.

\begin{table}[t]\centering\footnotesize
\caption{Pilot run on the 11 Windows tasks (Claude Opus 5.5, effort \emph{high}, full configuration). Quality is 0 for a wrong verdict, so every verdict here was correct.}
\label{tab:pilot}
\resizebox{\columnwidth}{!}{%
\begin{tabular}{@{}llrrrr@{}}\toprule
Task & Expected & Quality & Effic. & Reward & Calls\\\midrule
\texttt{logon-script-persistence} & malicious & 0.917 & 0.714 & 0.864 & 7\\
\texttt{lsass-dump-comsvcs} & malicious & 1.000 & 0.714 & 0.943 & 7\\
\texttt{service-binpath-hijack} & malicious & 0.922 & 0.600 & 0.849 & 10\\
\texttt{hta-startup-folder} & malicious & 0.922 & 0.500 & 0.830 & 12\\
\texttt{firewall-rule-added} & malicious & 0.887 & 0.667 & 0.828 & 6\\
\texttt{bits-download} & malicious & 0.929 & 0.800 & 0.891 & 5\\
\texttt{run-key-and-beacon} & malicious & 0.908 & 0.625 & 0.840 & 8\\
\texttt{lsass-query-by-svchost} & benign & 1.000 & 1.000 & 1.000 & 3\\
\texttt{python-http-server} & malicious & 0.922 & 0.455 & 0.822 & 11\\
\texttt{w32time-runtime-key} & benign & 1.000 & 0.750 & 0.950 & 4\\
\texttt{gpsvc-service-host} & benign & 0.850 & 0.429 & 0.753 & 7\\
\midrule
Mean & & 0.933 & 0.659 & 0.870 & 7.27\\
\bottomrule\end{tabular}}\end{table}


\paragraph{Where reward is lost (RQ2)} Since $R=Q(0.8+0.2E)$, the gap between the mean reward and its maximum
of 1 decomposes into a quality part, $\overline{1-Q}=0.068$, and an efficiency part,
$\overline{0.2\,Q(1-E)}=0.063$. Efficiency therefore costs almost as much as imperfect artifacts.
\ifconf The per-task counts show why\else Fig.~\ref{fig:pilotcalls} shows why\fi: the agent used 80 tool calls against a target of 49 (the sum of
$n^\star+1$), and 10 of 11 tasks exceeded their target. The largest overruns were the HTA execution (12 calls
against a target of 6), the Python HTTP server (11 against 5) and the service-configuration change (10 against
6), all multi-step chains in which the agent fetched relations one at a time. Benign tasks were cheaper
(4.7 calls on average, efficiency 0.73) than malicious ones (8.3 calls, efficiency 0.63), except the Group
Policy service host (7 calls against a target of 3), where the agent kept collecting evidence after the parent
process alone settled the question.

\paragraph{Quality} Quality was 1.0 on three tasks and at least 0.85 on all. With a correct verdict, the
score leaves only two ways to lose quality: for malicious tasks, missing core entities or naming entities
outside the acceptable set; for benign tasks, citing only some anchors or non-decisive fields. The per-call
traces of the study runs attribute these losses per task.

\ifconf\else
% Pilot: tool calls per task against the optimal count (pgfplots, reads data/pilot_per_task.csv).
\begin{figure}[t]
\centering
\begin{tikzpicture}
\begin{axis}[
  xbar, bar width=4pt, width=0.62\columnwidth, height=7.2cm,
  xmin=0, xmax=13, xlabel={Tool calls}, xmajorgrids, grid style={gray!25},
  symbolic y coords={gpsvc-service-host,w32time-runtime-key,lsass-query-by-svchost,python-http-server,run-key-and-beacon,bits-download,firewall-rule-added,hta-startup-folder,service-binpath-hijack,lsass-dump-comsvcs,logon-script-persistence},
  ytick=data, yticklabel style={font=\tiny\ttfamily}, tick label style={font=\scriptsize},
  legend style={font=\scriptsize, at={(0.98,0.02)}, anchor=south east, draw=gray!40},
  enlarge y limits=0.06,
]
\addplot[fill=cA, draw=none] table[col sep=comma, x=calls, y=task] {data/pilot_per_task.csv};
\addplot[fill=cB, draw=none] table[col sep=comma, x expr=\thisrow{optimal}+1, y=task] {data/pilot_per_task.csv};
\legend{Used, Optimal + 1}
\end{axis}
\end{tikzpicture}
\caption{Pilot run: tool calls used per task against the efficiency target (optimal + 1, at which efficiency
reaches 1). Every verdict was correct; the gap between the bars is the efficiency loss.}
\label{fig:pilotcalls}
\end{figure}

\fi
